\chapter{Эпистемика: уверенность, провенанс и ревизия}
\label{ch:epistemics}

\section{Ложные правдоподобные убеждения}

Живой ум держит массу ошибочных, но правдоподобных убеждений. Классический пример:
расхожее мнение, будто некий знаменитый предприниматель был «просто перформером, а
не инженером», — миф, рассыпающийся при чтении биографии. Стремиться к памяти,
вовсе не порождающей таких убеждений, бессмысленно и вредно: это было бы стерильно.
Задача — не запретить их, а научиться выставлять уверенность меньше единицы и потом
проверять.

\section{Недоопределённость и преждевременный коллапс}

Пусть правильный вывод — функция нескольких посылок:
\[
A \;=\; \text{logical\_f}(B, C, D, E),
\]
но уму известны только $B, C$, и он делает вывод $X = \text{logical\_f}(B, C)$.
Беда в том, что $B, C$ совместимы не только с $A$, но и с $X, Y$. Вывод
\emph{недоопределён} известными посылками (\emph{underdetermination}; философское
ядро — тезис Дюэма–Куайна). Когнитивное проявление — \emph{преждевременный коллапс}:
ум схлопывает суперпозицию совместимых гипотез в одну — часто удобную — слишком
рано. Формальное ядро: \emph{совместимость $\ne$ подтверждение} (родственник ошибки
«утверждение консеквента»).

\begin{pitfall}[Известная и неизвестная недоопределённость]
\textbf{Известная}: мы \emph{знаем}, что $D,E$ существуют, но не знаем их значений;
тогда можно честно выставить $\conf(X)$ низкой и пометить
$\texttt{need\_to\_discover}(D,E)$. \textbf{Неизвестная} — злой случай: уму
\emph{вообще неизвестно}, что $D,E$ существуют; в поле информации только $B,C$;
пометить нечего, уверенность может быть высокой. Поэтому per-inference калибровка
\emph{не} решает проблему целиком.
\end{pitfall}

\begin{workedbox}[Миф о предпринимателе, по шагам]
$B$: «известен маркетингом», $C$: «богат». Вывод $X$: «получил всё благодаря
происхождению/пиару». Но те же $B,C$ совместимы и с $A$: «сильный инженер, к тому
же умеющий продавать». Различающие посылки $D$: «первая лаборатория — в родительском
гараже», $E$: «глубокое STEM-образование» — уму \emph{неизвестны}, он даже не
подозревает об их существовании и уверенно выбирает удобный $X$. Ни низкая
$\conf(X)$, ни $\texttt{need\_to\_discover}$ здесь не помогут: не на что их повесить.
Спасает лишь квартет ниже.
\end{workedbox}

\section{Квартет починки}

Неизвестную недоопределённость нельзя закрыть в момент коммита; её лечат
статистически и постфактум.

\begin{enumerate}
\item \textbf{Prior скромности на класс вывода.} Ни один абдуктивный/индуктивный
коммит не получает уверенность $1$, как бы хорошо ни сидели данные:
$\conf(X)\le 1-\epsilon$ для всего \emph{сорта} «вывод по неполным посылкам».
«Согласуется, и альтернатив не придумывается» — ровно то состояние, что
предшествует «а потом выяснилось». Не требует знать $D,E$.
\item \textbf{Ревизуемость и понижение по противоречию.} Коммит делаем понижаемым:
позднее наблюдение, столкнувшееся с $X$, ретроактивно вскрывает недоопределённость,
и машинерия ревизии понижает узел. Починка ниже по потоку, не на записи.
\item \textbf{Телепорт-скраб.} Единственный поиск, дотягивающийся до
\emph{непредставленных} альтернатив (глава~\ref{ch:topology}).
\item \textbf{Внешний вход.} Кто-то знает $D,E$ и говорит; любопытство находит это
во внешнем источнике.
\end{enumerate}
Первый механизм держит \emph{известную} неопределённость; вместе с остальными тремя
— и \emph{неизвестную}. Полный разбор этой петли на живом примере (гипотеза
рождается, проверкой усиливается, затем опровергается) — «House\,/\,Cast~Away» в
главе~\ref{ch:fusion}.

\section{Система поддержки истинности и «ползучая вера»}

Отдельная беда — \emph{устойчивость убеждения к опровержению его основания}
(belief perseverance): когда неверный вывод $X$ уже продиффундировал в память и дал
потомков, эти потомки продолжают транслироваться даже после того, как $X$ признан
ложным.

Механически это \emph{проблема ретракции}, и для неё есть готовый аппарат —
\emph{система поддержки истинности} (TMS, Doyle; ATMS, de~Kleer). При выводе
следствия записываем его \emph{обоснование} — множество узлов, из которых оно
получено (расширенный провенанс, раздел~\ref{sec:prov}). Тогда ретракция посылки
$X$ автоматически снимает опору с её обоснованных потомков.

\begin{lstlisting}[style=pseudo]
def retract(G, X):
    demote(G, X)
    for c in G.justificands[X]:          # кто опирался на X
        if not still_supported(G, c):    # осталось ли иное обоснование
            mark(G, c, "need_to_revise")
            retract(G, c)                # рекурсивно вниз по обоснованиям
\end{lstlisting}

\begin{intuition}[Полный TMS слишком «нечеловечен»]
Полная система поддержки истинности мгновенно перестраивает всё мировоззрение при
опровержении любой посылки. Живой ум так не делает — и потому у него есть ползучая
вера. Значит нужен \emph{намеренно дырявый} TMS: обоснования пишем для важных и
высокоуверенных выводов, тривиальные не трекаем, а нетрекнутых потомков лечим
фоном.
\end{intuition}

Отсюда две «лечилки» разной латентности:
\begin{itemize}
\item \textbf{Retrieval-time augmentation} — патч \emph{на выходе}: при каждом
извлечении подмешивать свежую поправку «кстати, недавно я узнала $\neg X$». Дёшево,
чинит текущий разговор, хранилище не трогает.
\item \textbf{Ночной скраб} — ремонт \emph{хранилища}: медленно, случайными
инспекциями находит нетрекнутых потомков и понижает их.
\end{itemize}
Обе нужны.

\section{Эстафета к главе о личности}

Калиброванная уверенность и компульсивная перепроверка — не только эпистемическая
гигиена, но и \emph{черта характера}, выводимая из мотивации: «не „я уверена“, а
„похоже на $X$, но $Y$ я не проверила — дай гляну“». В главе~\ref{ch:personality}
мы увидим, что эпистемическая аккуратность Рин должна \emph{проявляться} в её
голосе, а не только жить во внутренней машинерии.
